\section{Tritic generators, quadratic actions, and energy transport}
\label{sec:hamiltonianos-curvatura}

Dynamic geometry preserves the distinction between the state generating a
realization and the coordinates in which that realization expresses its evolution. APP
constructs and preserves the additive and multiplicative sheets; the TRIT orients the
operating mode; the TPK composes selection, transport, and update on
the enriched state through
\[
 U_t=\operatorname{Upd}_t\circ\operatorname{Tra}_t\circ
 \operatorname{Sel}_t.
\]
Nonadic holonomy extends this state through
\(w_6\to w_{12}\to w_{18}\to w_{24}\to w_{30}\to R_{36}\to G_9\).
Phase return preserves the advance of memory. Its outputs
\(\pi_{\rm HMT},\varphi_{\rm HMT},e_{\rm HMT},\alpha_{\rm HMT}\)
belong to the same correlated generation, with the readers and internal
constructive order already established. The joint discrete structure
of the continuum subsequently supports the geometric and physical representations.
At this stage of the construction, a Hamiltonian is the functional
realizing an already typed transport on a symplectic fiber; its domain,
orientation, and time parameter are part of the statement.

The article on the holographic extreme and mean ratio recovers three
concrete operators. Multiplication by \(4\) modulo \(9\) acts on
the ternary orbits of units; on its augmentation module, its matrix is
\[
 C=\begin{pmatrix}0&-1\\1&-1\end{pmatrix},\qquad C^2+C+I=0.
\]
Parabolic transport and areal--volumetric incidence are
represented, on their respective planes, by
\[
 N=\begin{pmatrix}0&1\\0&0\end{pmatrix},\qquad
 F_{\rm av}=\begin{pmatrix}0&1\\1&1\end{pmatrix}.
\]
Incidence originates in the sheet calendar; its powers count
compositions of these modes. The resulting real generators are
\begin{equation}
 J_+=\frac{2C+I}{\sqrt3},\qquad J_0=N,\qquad
 J_-=\frac{2F_{\rm av}^{2}-3I}{\sqrt5},\qquad
 J_\tau^2=-\tau I.
 \label{eq:hc-generadores}
\end{equation}
Here \(\tau=+1,0,-1\) distinguishes the elliptic, parabolic, and
hyperbolic regimes. The fibers have their own origins and bases. The formulas
\(e^{\pi J_+}=-I\), \(e^{J_0}=I+J_0\), and
\(e^{2\log\varphi J_-}=F_{\rm av}^{2}\) evaluate previously generated HMT
coordinates on these operators~\cite{hmtX}.

\subsection{Alternating form and Hamiltonian of the generator}

On an oriented two-dimensional real fiber, fix
\[
 \omega=dQ\wedge dP,\qquad
 \Omega=\begin{pmatrix}0&1\\-1&0\end{pmatrix},\qquad z=(Q,P)^T.
\]
We adopt \(\iota_{X_H}\omega=dH\); therefore
\(\dot Q=\partial_PH\), \(\dot P=-\partial_QH\).
This convention also fixes the sign of the oriented action.

\begin{proposition}[Hamiltonian realization of a traceless generator]
\label{prop:hc-cuadratica}
For a real two-by-two traceless matrix \(J\), the matrix
\(G_J=-\Omega J\) is symmetric, and the Hamiltonian
\begin{equation}
 H_J(z)=\frac12z^TG_Jz
 \label{eq:hc-hj}
\end{equation}
produces exactly \(\dot z=Jz\). If \(J^2=-\tau I\),
then \(\det G_J=\tau\). For a symplectic isomorphism
\(B:V\to V'\), the transport
\[
 J'=BJB^{-1},\qquad H'(z')=H_J(B^{-1}z')
\]
satisfies \(G'=B^{-T}G_JB^{-1}\) and intertwines both flows.
\end{proposition}
\begin{proof}
Writing \(J=\bigl(\begin{smallmatrix}a&b\\c&-a\end{smallmatrix}\bigr)\)
gives \(G_J=\bigl(\begin{smallmatrix}-c&a\\a&b\end{smallmatrix}\bigr)\).
Thus \(\Omega\nabla H_J=\Omega(-\Omega J)z=Jz\).
Cayley--Hamilton gives \(J^2=-\det(J)I\), and
\(\det(-\Omega)=1\). Finally, \(B^T\Omega B=\Omega\)
allows the substitution \(z=B^{-1}z'\) in both the form and the equation.
\end{proof}

Applied to \eqref{eq:hc-generadores}, the proposition gives
\begin{equation}
 H_{J_+}=-\frac{Q^2-QP+P^2}{\sqrt3},\qquad
 H_{J_0}=\frac{P^2}{2},\qquad
 H_{J_-}=\frac{-Q^2-QP+P^2}{\sqrt5}.
 \label{eq:hc-hamiltonianos-nativos}
\end{equation}
The elliptic sign records the orientation of \(C\) in the augmentation basis.
The reverse flow \(-J_+\) has the opposite positive form. The action
ellipse of the corpus uses precisely a decreasing phase orientation
for its positive Hamiltonian. Preserving this sign makes
the cyclic operator compatible with the action \(P\,dQ\).

Conjugation also provides an exact comparison criterion.
If \(BJ_\tau=J_{\tau'}B\) and \(B\) is invertible, squaring
gives \((\tau-\tau')B=0\), hence \(\tau=\tau'\). A connection may
preserve \(\omega\) between fibers in different regimes and preserve
their respective memories; conjugation of the generators additionally requires
equality of the regime. This distinction makes it possible to study a
tritic transition without replacing it by a renaming of the same flow.

\subsection{Deformation of the ellipse and Legendre transformation}

The coordinate \(\eta=\operatorname{artanh}(C^*/A)\), on the domain
\(|C^*/A|<1\), defines the symplectic frame
\(M_\eta=\operatorname{diag}(e^{\eta/2},e^{-\eta/2})\).
Deformation in the two directions is reciprocal and preserves their product.
In the normalized chart \((q,p)=M_\eta^{-1}(Q,P)\), consider
\[
 L_\tau=\begin{pmatrix}0&1\\-\tau&0\end{pmatrix},\qquad
 I_{\tau,\eta}=\frac12\bigl(\tau e^{-\eta}Q^2+e^\eta P^2\bigr).
\]
These representatives classify oriented quadratic flows. Their use
on a native fiber proceeds through a declared symplectic frame,
in accordance with Proposition~\ref{prop:hc-cuadratica}. The positive
elliptic representative is the one used by the action ellipse; for this representative, \(I\)
has the dimension of action.

The frame changes can be written explicitly. If
\(b_+=\sqrt{\sqrt3/2}\) and \(b_-=\sqrt{\sqrt5/2}\), the matrices
\[
 B_+=\begin{pmatrix}b_+&b_+/\sqrt3\\0&2b_+/\sqrt3\end{pmatrix},
 \qquad
 B_-=\begin{pmatrix}b_-&-b_-/\sqrt5\\0&2b_-/\sqrt5\end{pmatrix}
\]
have determinant one and satisfy
\( (-J_+)B_+=B_+L_{+1}\) and
\(J_-B_-=B_-L_{-1}\). For the parabolic regime, take
\(B_0=I\). Multiplication of the two columns verifies the
equalities, and determinant one proves preservation of the alternating
form.

\begin{theorem}[Variable transport and regular Lagrangian]
\label{thm:hc-legendre}
Let \(\eta(t)\) be of class \(C^1\) and \(\omega(t)>0\). The Hamiltonian
\begin{equation}
 H_{\tau,\rm tr}
 =\omega I_{\tau,\eta}+\frac{\dot\eta}{2}QP
 \label{eq:hc-htr}
\end{equation}
preserves \(I_{\tau,\eta(t)}\) along its solutions. Its
Legendre transformation with respect to \(P\) is invertible and gives
\begin{equation}
 \mathcal L_{\tau,\rm tr}(Q,\dot Q,t)
 =\frac{(\dot Q-\dot\eta Q/2)^2}{2\omega e^\eta}
 -\frac{\omega\tau e^{-\eta}}2Q^2.
 \label{eq:hc-ltr}
\end{equation}
\end{theorem}
\begin{proof}
The equations are
\[
 \dot Q=\omega e^\eta P+\frac{\dot\eta}2Q,
 \qquad
 \dot P=-\omega\tau e^{-\eta}Q-\frac{\dot\eta}2P.
\]
Differentiating \(I_{\tau,\eta(t)}\) and substituting these equations cancels separately
the \(\omega QP\) terms and the terms involving variation of \(\eta\).
The second derivative \(\partial_P^2H=\omega e^\eta>0\)
allows one to solve for
\(P=(\dot Q-\dot\eta Q/2)/(\omega e^\eta)\).
Substitution into \(P\dot Q-H\) proves \eqref{eq:hc-ltr}.
Moreover,
\[
 P\dot Q-H_{\tau,\rm tr}
 =p\dot q-\frac\omega2(\tau q^2+p^2),
\]
which proves that the mixed term is the variational connection of the moving
frame. Conservation follows from the complete transport, not from keeping
the semiaxes constant.
\end{proof}

The elliptic result and its action are developed in the variational
corpus~\cite{hmtX}; the same proof exhibits its two conjugate quadratic
realizations. The vanishing determinant of the parabolic form
coexists with a regular Legendre transformation: the latter depends on the
second derivative with respect to momentum, not on the determinant of the complete form.
By contrast, for the linear cotangent lift
\(H_A(q,p)=p^TAq\), one has \(\partial_p^2H_A=0\).
Its first-order action preserves \(\dot q=Aq\) through a
multiplier; it is a variational object distinct from
\eqref{eq:hc-ltr}.

\subsection{Memory return and conserved variational charge}

In the calendar memory record, the counters \((g,s)\)
receive the increment \((4,72)\) upon a complete return. The relation
\(72=18\cdot4\) fixes
\[
 v=(1,18)^T,\qquad I_{\rm mem}=s-18g,
 \qquad B=\begin{pmatrix}1&0\\18&1\end{pmatrix}.
\]
The native stabilizer is
\begin{equation}
 N_v=BNB^{-1}=
 \begin{pmatrix}-18&1\\-324&18\end{pmatrix},\qquad
 N_vm=vI_{\rm mem}(m),\qquad N_v^2=0.
 \label{eq:hc-nv}
\end{equation}
On the realified plane \(m=(g,s)^T\), equipped with \(dg\wedge ds\),
the Hamiltonian \(H_{\rm mem}=I_{\rm mem}^2/2\) produces
\(dm/du=N_vm\). The parameter \(u\) belongs to this dimensionless
realization of the record. The chart
\(q=g,\ p=s-18g\) is symplectic and transforms the action one-form
according to
\begin{equation}
 s\,dg=p\,dq+d(9q^2).
 \label{eq:hc-memoria-borde}
\end{equation}
Thus, eliminating \(s\) from the first-order action gives
\begin{equation}
 \mathcal L_{\rm mem}
 =\frac12\left(\frac{dg}{du}\right)^2
 +18g\frac{dg}{du}
 =\frac12\left(\frac{dg}{du}\right)^2
 +\frac{d}{du}(9g^2).
 \label{eq:hc-lmem}
\end{equation}
The variational equation is \(d^2g/du^2=0\), and the momentum
conserved in the symplectic chart is \(p=I_{\rm mem}\).
The complete affine return originates in the linear Hamiltonian
\(H_{\rm ret}=4I_{\rm mem}\): its time-\(u=1\) map
is \((g,s)\mapsto(g+4,s+72)\). Both Hamiltonians commute
because they are functions of the same charge, but they realize distinct
operations: parabolic stabilization and translation of the record.
These formulas recover the variational mechanism of memory, including
its boundary term~\cite{hmtX}.

\subsection{Metric curvature and geodesic energy}

The metric realization of the regimes uses a radial chart and an
inverse length scale \(\rho>0\):
\begin{equation}
 g_{\tau,\rho}=dr^2+S_{\tau,\rho}(r)^2d\theta^2,
 \qquad
 S_{\tau,\rho}(r)=
 \begin{cases}
 \sin(\rho r)/\rho,&\tau=+1,\\
 r,&\tau=0,\\
 \sinh(\rho r)/\rho,&\tau=-1.
 \end{cases}
 \label{eq:hc-metrique}
\end{equation}
The Gaussian curvature is \(K_G=-S''/S=\tau\rho^2\).
The chart is considered in its regular interior, \(S>0\); for the
spherical realization, \(0<r<\pi/\rho\). The poles are treated using
different regular charts. The parameter \(\rho\) and the dimensional
realization of the metric retain their own specification.

For a realized mass \(m>0\), the geodesic Lagrangian and its
Hamiltonian are
\begin{equation}
 \mathcal L_{\rm geo}=\frac m2(\dot r^2+S^2\dot\theta^2),\qquad
 H_{\rm geo}=\frac1{2m}\left(p_r^2+\frac{p_\theta^2}{S^2}\right).
 \label{eq:hc-geodesique}
\end{equation}
Indeed, \(p_r=m\dot r\), \(p_\theta=mS^2\dot\theta\);
the velocity Hessian is \(m\operatorname{diag}(1,S^2)\)
and is positive and invertible. The radial equation is
\(\ddot r-SS'\dot\theta^2=0\), whereas
\(d(mS^2\dot\theta)/dt=0\) conserves angular momentum.
Positivity of this energy holds for all three curvatures. The indefinite
form of the hyperbolic generator in \eqref{eq:hc-hamiltonianos-nativos}
and the positive geodesic energy of a negatively curved metric
are functions on different spaces. This comparison determines precisely
which meaning of curvature enters each statement.

\subsection{Constitutive response and modal realization of Maxwell theory}

The action sections and angular coordinates precede the
constitutive response \cite{hmtIntegral}. In the fixed real orientation \(x>y>0\),
with the projectors of the two sheets,
\[
 T=q_+P_++q_-P_-,\quad q_\pm=e^{-x\mp y},\quad
 \mathcal V=(I-T^{90})(I-T^{120})^{-1},
 \qquad 0<q_+<q_-<1,
\]
the functional calculus produces
\[
 r_\pm=\frac{1-q_\pm^{90}}{1-q_\pm^{120}},\qquad
 \widehat\varepsilon=r_+^2,\quad \widehat\mu=r_-^2,
 \quad \widehat c=(r_+r_-)^{-1},\quad \widehat Z=r_-/r_+.
\]
The corresponding dimensional charts preserve exactly
\(\varepsilon\mu=c^{-2}\) and \(\mu/\varepsilon=Z^2\).
Permittivity, permeability, and velocity are thus outputs
of the same operator, with the incidence sectors fixed by the TPK;
the following dynamics uses them as already constructed results.

To exhibit transport to a field Hamiltonian, fix an
oriented spatial realization, boundary conditions allowing
integration by parts, and a real transverse mode \(f\) normalized by
\[
 \nabla\cdot f=0,\qquad \int|f|^2=1,\qquad
 \nabla\times\nabla\times f=k^2f,\qquad k>0.
\]
In transverse gauge and without sources, \(A=af\) and the momentum
\(\Pi=pf\) define an immersion of the modal plane into field
space. For homogeneous constitutive constants, the Maxwell Hamiltonian and
symplectic form restrict to
\begin{equation}
 H_{\rm EM}(a,p)=\frac{p^2}{2\varepsilon}
 +\frac{k^2a^2}{2\mu},\qquad \omega_{\rm EM}=da\wedge dp.
 \label{eq:hc-maxwell-mode}
\end{equation}
Integration by parts gives
\(\int|\nabla\times f|^2=k^2\), proving this restriction.
With \(\omega_k=ck\), the change
\begin{equation}
 Q=\sqrt{\varepsilon\omega_k}\,a,
 \qquad P=\frac{p}{\sqrt{\varepsilon\omega_k}}
 \label{eq:hc-maxwell-symplectic}
\end{equation}
preserves \(da\wedge dp=dQ\wedge dP\) and gives
\(H_{\rm EM}=\omega_k(Q^2+P^2)/2\). Its inverse and its equations
\(\dot Q=\omega_kP,\ \dot P=-\omega_kQ\)
constitute the modal intertwining with the positive elliptic regime.
The restriction is invariant because the Maxwell operator preserves the
transverse eigenmode. The Legendre transformation produces
\(\mathcal L_{\rm EM}=\varepsilon\dot a^2/2-k^2a^2/(2\mu)\).

Two transverse modes of equal frequency, combined in quadrature,
realize the two circular orientations of the corpus. The field satisfies
\(\mathbf B=c^{-1}\widehat{\mathbf k}\times\mathbf E\) and
\(\omega=ck\). The relation to the oscillator is proved here through
the field immersion, modal normalization, and transport of the symplectic
form, with the boundary conditions kept explicit.
A change in \(k\), cavity geometry, or preparation
modifies the realized mode; it leaves the genealogy of
\(\varepsilon,\mu,c,Z\) intact.

In the Lorentzian chart of the same sector, coupling to a material
section is written using \(F=dA\) and an already typed charge \(q_e\).
For this variational realization, \(\nabla\) is taken to be the
Levi--Civita connection of \(g\). On
\(T^*\mathcal M\), the covariant Hamiltonian
\[
 H_A=\frac1{2m}g^{\mu\nu}(p_\mu-q_eA_\mu)(p_\nu-q_eA_\nu)
\]
is the Legendre transform of
\(\mathcal L_A=\frac m2g_{\mu\nu}\dot x^\mu\dot x^\nu
+q_eA_\mu\dot x^\mu\). The Hessian \(mg\) is invertible.
The Euler--Lagrange equation gives
\(m\nabla_u u^\mu=q_eF^\mu{}_{\nu}u^\nu\).
Antisymmetry of \(F\) preserves \(g(u,u)\). This is the
covariant realization with an affine parameter, whereas
\eqref{eq:hc-maxwell-mode} uses the time of a transverse spatial
decomposition: both parameters and both domains are defined.

\subsection{Graded channels, Clifford representation, and mass operator}

The material atlas first produces complete channels of route, sheet, flavor,
color, and representation. At finite depth, its Hermitian
realization is \(\mathcal H_K=\bigoplus_{s\in S_K}V_s\), with an
orthonormal channel basis and transported parity grading. The exterior
algebra of the odd sector generates the CAR: insertion of a mode
\(a_j^*\) and its adjoint removal \(a_j\) satisfy
\(\{a_i,a_j^*\}=\delta_{ij}I\) and \(\{a_i,a_j\}=0\).
The proof counts the exchange signs of complete modes; the masses
are evaluated subsequently~\cite{hmtVI}.

In the single-mode exterior sector, the basis \((|0\rangle,|1\rangle)\)
identifies the removal operator with
\(a=N\). The map \(e_1\mapsto|0\rangle,\ e_2\mapsto|1\rangle\)
is an isometry of the normalized complex plane and intertwines both
nilpotent operators. This identification uses the declared inner product and
grading; it preserves the distinct meanings of their bases.
With the internal \(i\) already realized, one obtains
\begin{equation}
 \sigma_1=N+N^*,\qquad
 \sigma_2=-i(N-N^*),\qquad
 \sigma_3=NN^*-N^*N.
 \label{eq:hc-pauli}
\end{equation}
Direct multiplication gives
\(\sigma_j^*=\sigma_j\) and
\(\sigma_i\sigma_j=\delta_{ij}I+i\epsilon_{ijk}\sigma_k\).
In particular, the nilpotent pair and its adjoint produce a
Clifford representation; the adjoint is that of the Hermitian fiber and
is transported together with it.

The mass operator on a finite collection of channels \(\mathcal F\)
is constructed by the record, character, and their refinements:
\begin{equation}
 M_\beta=S M_0 S,\qquad
 S=R_{\rm act}^{B_M/2},\quad
 R_{\rm act}=\hbar_{\rm pre}/\hbar_{\rm ret}>0,
 \qquad M_0>0.
 \label{eq:hc-mass}
\end{equation}
Here \(B_M\) is the diagonal route-exponent operator, and \(M_0\)
contains the refined character relative to the electronic origin.
Positivity follows from \(\langle v,M_\beta v\rangle
=\langle Sv,M_0Sv\rangle>0\) for \(v\ne0\).
The sheet product preserves the null sector through its own realization;
the exponential of the material character remains positive.

In this realization, write \(\hbar:=\hbar_{\rm ret}\),
the already constructed action section, and \(c:=c_{\rm vac}\) in the
declared dimensional chart. On \(\mathbb C^4\otimes\mathcal F\), define
\[
 \alpha_j^{D}=\sigma_1\otimes\sigma_j,\qquad
 \beta^D=\sigma_3\otimes I_2.
\]
The superscript distinguishes these matrices from the fine-structure constant.
The spatial chart \(1+3\) assigns three momentum components
\(\mathbf p\). The multiplication map
\begin{equation}
 H_D(\mathbf p)
 =c\sum_{j=1}^3\alpha_j^Dp_j\otimes I_{\mathcal F}
 +c^2\beta^D\otimes M_\beta
 \label{eq:hc-dirac}
\end{equation}
is the free Dirac realization of this material operator.

\begin{theorem}[Spinorial factorization of material dispersion]
\label{thm:hc-dirac-square}
The matrix \eqref{eq:hc-dirac} is self-adjoint and satisfies
\begin{equation}
 H_D(\mathbf p)^2
 =I_4\otimes\bigl(c^2|\mathbf p|^2I_{\mathcal F}
 +c^4M_\beta^2\bigr).
 \label{eq:hc-dirac-square}
\end{equation}
On an eigenspace \(M_\beta v=m v\), its energies are
\(\pm E_m(\mathbf p)\), where
\(E_m=\sqrt{c^2|\mathbf p|^2+m^2c^4}\).
\end{theorem}
\begin{proof}
The relations in \eqref{eq:hc-pauli} give
\(\{\alpha_i^D,\alpha_j^D\}=2\delta_{ij}I_4\),
\(\{\alpha_i^D,\beta^D\}=0\), and \((\beta^D)^2=I_4\).
The flavor factor commutes with the spinorial matrices.
The cross terms vanish, and the squares give
\eqref{eq:hc-dirac-square}. The vanishing trace and Clifford relations
produce both branches, with spinorial multiplicity two when \(E_m>0\).
\end{proof}

In momentum representation, the natural domain is
\(\{\psi\in L^2(\mathbb R^3;\mathbb C^4\otimes\mathcal F):
H_D\psi\in L^2\}\). As a measurable Hermitian matrix acting
by multiplication, the operator is self-adjoint on this domain. The
spatial representation originates in the unitary Fourier transform
subsequent to the construction of the channels. The theorem determines the free
realization and its exact dependence on HMT mass; the interaction retains its
connection operators and its own domain.

On the positive branch \(m>0\), the ordinary Legendre transformation
of the symbol \(E_m(\mathbf p)\) is
\[
 \mathbf v=\frac{c^2\mathbf p}{E_m},\qquad
 \mathcal L_m=-mc^2\sqrt{1-|\mathbf v|^2/c^2},
 \qquad |\mathbf v|<c.
\]
Solving \(\mathbf p=m\mathbf v/\sqrt{1-|\mathbf v|^2/c^2}\)
and substituting into \(\mathbf p\cdot\mathbf v-E_m\) proves the
equality. The spinorial action, by contrast, is first order in the field:
\(\mathcal L_D=\frac{i\hbar}{2}(\psi^*\dot\psi-\dot\psi^*\psi)
-\psi^*H_D\psi\). Its variation gives
\(i\hbar\dot\psi=H_D\psi\); its Hessian in the field velocities
is singular. Both statements are compatible because they act on
different variational spaces.

For a static Abelian connection in a flat spatial chart, with
\(\Pi_j=-i\hbar\partial_j-q_eA_j\), constant mass, and vanishing
scalar potential, the following holds on smooth compactly supported functions:
\[
 [\Pi_i,\Pi_j]=iq_e\hbar\epsilon_{ijk}B_k,
\]
and the same factorization produces
\begin{equation}
 H_D(A)^2=c^2\boldsymbol\Pi^2+c^4M_\beta^2
 -q_e\hbar c^2\boldsymbol\Sigma\cdot\mathbf B,
 \qquad \Sigma_j=I_2\otimes\sigma_j.
 \label{eq:hc-dirac-curvature}
\end{equation}
The flavor identity factors are implicit. The curvature term
is the product of two antisymmetric objects: the commutator of covariant
derivatives and Pauli multiplication. Its coefficient is fixed
when the charge representation and action chart are chosen; it is not
selected through a target spectral value.

If \(V\) is the already constructed unitary flavor transformation, then
\[
 (I_4\otimes V)H_D(M_\beta)(I_4\otimes V^*)
 =H_D(VM_\beta V^*).
\]
Unitary mixing preserves the complete spectrum. The positive congruence
in \eqref{eq:hc-mass} transports a different structure and may modify the
eigenvalues in a fixed metric. For example,
\(M_0=\operatorname{diag}(1,2)\), \(S=\operatorname{diag}(2,1)\)
gives \(SM_0S=\operatorname{diag}(4,2)\). If the inner product
is also transported to \(G'=S^*S\), the generalized problem
\(SM_0Sv=\lambda G'v\) recovers the eigenvalues of \(M_0\)
through \(w=Sv\). In this way, a change of flavor basis, material
deformation, and metric transport have precise spectral consequences.

\subsection{Curvature of the scalar potential in the electroweak realization}

The electroweak sector uses the determinant of the same angular operator,
\(D_A=\det T=e^{-2A_{\rm rad}}\), together with the already constructed
material signatures. In the angular realization with common refinements of
\(W\) and \(Z\), their signatures \((94,-1)\) and \((95,-1)\) give
\(m_W^\angle/m_Z^\angle=e^{-A_{\rm rad}}\).
In the rationalized tree-level convention of the corpus,
\[
 g^2=\frac{4\pi_{\rm HMT}\alpha_{\rm HMT}}{1-D_A},\qquad
 \lambda_H=\frac{\pi_{\rm HMT}\alpha_{\rm HMT}}{2(1-D_A)}
 \exp\!\left(6A_{\rm rad}+\frac{10}3C^*_{\rm rad}\right)>0.
\]
The second formula follows from the scalar signature \((97,9)\) and
\(m_H^2/m_W^2=8\lambda_H/g^2\). The potential normalization is
\(V(\mathsf H)=-\mu_H^2\mathsf H^*\mathsf H
+\lambda_H(\mathsf H^*\mathsf H)^2\). These compositions preserve
their scale and tree-level convention; the refinements that cancel
in a ratio are declared in the realization itself.

In a real radial chart \(\mathsf H^*\mathsf H=h^2/2\), the potential
is \(V(h)=-\mu_H^2h^2/2+\lambda_Hh^4/4\). For
\(\mu_H^2>0\), its critical points are \(0\) and
\(\pm v\), with \(v^2=\mu_H^2/\lambda_H\), and
\[
 V''(0)=-\mu_H^2,\qquad V''(\pm v)=2\mu_H^2.
\]
In the natural units of this chart, the normalized homogeneous mode has
\(H_h=p_h^2/2+V(h)\). Its linearization at a critical point \(h_0\)
has generator
\[
 A_{h_0}=\begin{pmatrix}0&1\\-V''(h_0)&0\end{pmatrix},
 \qquad A_{h_0}^2=-V''(h_0)I.
\]
Thus the origin is hyperbolic and the minima are elliptic; the
boundary \(V''=0\) is of parabolic type. This statement follows
from the second derivative of the potential and an explicit modal
normalization. The positive material energy of the minimum and the unstable
direction of another background thus belong to the same function with different
expansion points. The linearized mode is a particular transport of
the three quadratic types; its frequencies and units originate in the
potential, not in identifying its Hessian with the Gaussian curvature in
\eqref{eq:hc-metrique}.

\subsection{Yang--Mills curvature mode and elliptic oscillation}

The Yang--Mills realization constructed in the preceding sections
has a positive generator \(D^{\rm YM}\) of inverse-length
dimension. The scale
\[
 \kappa_{\rm av}=\frac{\mu_{\rm vol}}{\ell_0}
 \left(\frac{\pi_{\rm HMT}}{\varphi_{\rm HMT}^2}-1\right)>0,
 \qquad \delta=3\kappa_{\rm av}
\]
descends from sheet incidence and the regional reader, with the declared
length unit \(\ell_0\). The vacuum is simple, and the curvature
witness in the product chart is
\[
 W(q)=\mathbf1_{\{q_1+\cdots+q_6=0\}}-\frac13,
 \quad \mathbb EW=0,\quad \|W\|^2=\frac29,
 \quad D^{\rm YM}W=\delta W.
\]
Unitary transport to the interaction measure takes \(W\) to
\(W_g\) and preserves these identities. The refinement inclusions
\(J_m\) intertwine the generators and preserve the witness, so
it persists in the physical limit. Here the complete result of the
vacuum, uniform-gap, and reconstruction proofs is used, rather than an
isolated eigenvalue of a finite matrix~\cite{HMT-YM}.

Let \(w=W_g/\sqrt{2/9}\) be the unit vector in the physical fiber and
\(\hbar=\hbar_{\rm ret}>0\). The energy Hamiltonian is
\(H_E=\hbar cD^{\rm YM}\). On the realified Hilbert space,
adopt the form \(\omega_{\mathcal H}(u,v)=2\hbar
\operatorname{Im}\langle u,v\rangle\), with an inner product antilinear
in the first argument. Define
\begin{equation}
 \mathfrak F(Q,P)=\frac{Q+iP}{\sqrt{2\hbar}}\,w.
 \label{eq:hc-ym-embedding}
\end{equation}

\begin{theorem}[Symplectic immersion of the physical curvature mode]
\label{thm:hc-ym-oscillator}
The map \(\mathfrak F:\mathbb R^2\to\mathcal H^{\rm phys}\)
is a real immersion whose image is invariant under physical evolution and
satisfies
\begin{equation}
 \mathfrak F^*\omega_{\mathcal H}=dQ\wedge dP,
 \qquad
 \langle\mathfrak F,H_E\mathfrak F\rangle
 =\frac{c\delta}{2}(Q^2+P^2).
 \label{eq:hc-ym-energy-pullback}
\end{equation}
The restricted Schrödinger equation exactly intertwines the elliptic
flow \(\dot Q=c\delta P,\ \dot P=-c\delta Q\).
The map is compatible with the refinement inclusions.
\end{theorem}
\begin{proof}
For two variations \((a,b),(a',b')\), the inner product of their
images has imaginary part \((ab'-ba')/(2\hbar)\).
This proves the first identity. The eigenvalue equation and \(\|w\|=1\)
give the second. Since
\(\dot\psi=-icD^{\rm YM}\psi\), the coefficient \(Q+iP\)
satisfies \(\dot Q+i\dot P=-ic\delta(Q+iP)\).
Finally, \(J_mw_m=w_{m+1}\) implies
\(J_m\mathfrak F_m=\mathfrak F_{m+1}\). These identities pass to the
limit through the already constructed isometries.
\end{proof}

The complete plane describes amplitudes of one mode; its circle
\(Q^2+P^2=2\hbar\) represents the normalized vectors of this
complex eigenspace. On this circle, the global phase evolves and the
quantum ray remains fixed. Superpositions of distinct eigenspaces
also retain their relative phases. The immersion preserves this
distinction between amplitude space and projective space.

The energy associated with the spectral gap and the mass gap are
\(\Delta_E=\hbar c\delta\) and \(m_{\rm gap}=\hbar\delta/c\).
The modal frequency \(c\delta\) has inverse-time dimension.
The Euclidean lengths in \(e^{-sD^{\rm YM}}\) and the physical time
in \(e^{-itH_E/\hbar}\) are related through the chart and
reconstruction, preserving their distinct groups and semigroups.
The composition \(\mathfrak F\circ M_\eta^{-1}\) transports this
same energy to the action ellipse; if the frame varies, the connection
term in \eqref{eq:hc-htr} completes its evolution.

The resulting connection is an exact realization on the invariant curvature
mode and its extensions. Positivity of the quantum Hamiltonian
generates a rotation in its symplectic realification; the designation
of a spatial metric as hyperbolic pertains to the metric tensor, and the
unipotent transport of the record pertains to its memory. In finite
dimension, a group \(I+tN\) unitary for every \(t\in\mathbb R\), with
\(N^2=0\), forces \(N=0\): differentiating unitarity gives \(N^*=-N\),
and then \(N^*N=-N^2=0\). The classical parabolic flow preserves its
symplectic form and has a quantization on suitable spaces; this
identity only delimits its identification with a finite unitary matrix.

A collection of concrete transports is thereby articulated. Memory
produces a charge and a parabolic action; the ellipse frame produces
the variational connection and its Lagrangian; the constitutive responses
determine the Maxwell coefficients; the channel grading constructs
Clifford and allows factorization of the material operator; Yang--Mills
reconstruction supplies a symplectic modal immersion preserving exactly
its gap and refinement. Each identification has a map,
a preserved form, and an intertwined equation.
